\section{Materials and methods}
\subsection{Study scope and controlled implementation}
NetSy (A Network loads and emissions Synthesizer) represents freight activity
through synthetic demand, vehicle tours, network loading and environmental
accounting. This study examines a Dutch application using the MASS-GT input
case and external gateways \parencite{massgt2025}. The comparison concerns correspondence with an
existing model and sensitivity to alternative formulations; it does not assess
accuracy against independent traffic or emissions observations.

The implementation is frozen at NetSy v0.5.0 (commit \texttt{491b1d5}).
The retained source reference is commit \texttt{4f727c3}, with execution
repairs documented separately in the audit. Source model-logic corrections
are identified explicitly rather than counted as methodological improvements.
Case inputs are copied into an isolated, hashed snapshot. Five upstream
realizations use stage seeds 1--5 for agent, shipment, freight-tour and parcel
scheduling synthesis; the minimum establishment size is four employees.
Parcel demand and service matrices are fixed by the supplied inputs.
Within a realization, each alternative reuses the saved freight tours, parcel
demand and service matrices. This pairing controls upstream synthesis rather
than assuming that equal seed integers produce equal random draws across
algorithms. The modelled period follows the supplied case and matrices;
totals are not presented as independently observed national daily inventories.

\subsection{Parcel scheduling and network paths}
The corrected inherited parcel scheduler is compared with the available
Clarke--Wright savings implementation \parencite{clarke1964}. The inherited method combines proximity clustering with pairwise stop swaps;
the savings method combines capacity-constrained route merging with single-stop
relocation, retaining direct tours for full zone loads. The contrast therefore
compares the implemented scheduling packages, including their local improvement
steps. Both retain the supplied depot/customer
allocation, parcel quantities, vehicle capacities and service parameters,
including depot returns. The supplied tour capacity is 180 parcels. Scheduling is capacity-constrained;
computed arrival times are not evidence that explicit customer time-window
constraints have been enforced. There is no explicit finite fleet-size
constraint; tour counts are not simultaneous vehicle requirements. They may produce different tour counts and departure
times. The supplied link-speed table is constant across time intervals, so
changed departure times do not introduce an additional speed-profile treatment.
The experiment examines the integration and consequences of these established
heuristics in the case application.

All-or-nothing loading uses Dijkstra shortest paths with the supplied time and
distance costs and vehicle restrictions \parencite{dijkstra1959}. A terrain sensitivity adds an uphill
preference to the link cost: link kilometres multiplied by positive signed
grade (in percent) and a coefficient of 0, 0.05 or 0.20 in the case's cost units
per kilometre per percentage point. Downhill coefficients remain zero.
Existing replacement costs for restricted links retain their original
precedence. These coefficients bracket an explicitly uncalibrated preference;
they are neither estimated energy costs nor a COPERT-based route objective.
Tour construction continues to use the supplied skims: this experiment changes
network paths after scheduling, without an iterative rescheduling feedback.

\subsection{Terrain data supplement}
The original network contains 158,806 directed links and no grade attribute.
AHN4 terrain and surface rasters at native 5\,m resolution supplement the case.
Files are selected through the official national download portal, with
provider SHA-256 archive checksums retained. AHN describes these as ground and
surface products in Dutch national coordinates with heights relative to NAP;
the surface includes objects such as vegetation and buildings.\footnote{AHN
product description: \url{https://www.ahn.nl/dataroom}; official download portal:
\url{https://basisdata.nl/hwh-portal/download/index.html}, accessed 19 November
2026.}

At each physical network node, the elevation estimate is the median of a
$3\times3$ raster neighbourhood, requiring at least five valid cells.
An estimate is retained only when the terrain and surface medians differ by
at most 3\,m. Signed grade is the endpoint elevation difference divided by
horizontal link geometry length, expressed as a percentage. Direction follows
source node identifiers: reverse arcs can share the same coordinate order in
the network geometry. Accepted reverse-arc grades are checked for opposite signs.
Links shorter than 50\,m, gradients above 10\% in magnitude, missing endpoints,
connectors, ferries and links named as bridges, viaducts or tunnels receive a
flagged zero modelling fallback. This screening is incomplete: the result is
a terrain-informed proxy, not a validated road-deck gradient or complete
structure inventory. Raster acquisition and transport-input dates need not
coincide. Coverage and fallback exposure are reported alongside effects.

\subsection{Emissions and comparison boundaries}
Building on integrated freight-network and COPERT assessment approaches
\parencite{dias2026}, each saved assignment is evaluated with both the inherited empty/full,
road-type factors and the supplied COPERT V hot-exhaust configuration. This follows the established
COPERT modelling family \parencite{eeaRoad2026}; the exact coefficient table and
vehicle mapping used here are retained with the hashed case inputs. The case
adapter extracts the July 2026 road-transport Appendix 4 workbook (COPERT 5.10).
Its CO$_2$ factors are derived from the energy-consumption functions using
74.1\,g CO$_2$/MJ for diesel \parencite{ipccMobile2006}.
The inherited method interpolates factors by utilization. COPERT uses supplied
vehicle mappings, speed-dependent functions, and nearest grade and payload
bins (grade: $-6,-4,-2,0,2,4,6$\%; payload: 0, 50, 100\%). Speeds outside
coefficient validity limits are clipped as specified in the implementation.
The supplied truck coefficients vary with grade. The van coefficients are
repeated across grade bins, so parcel emissions have no direct grade-factor
response in this configuration; path and speed changes can still affect them.
The supplied COPERT fleet assumptions represent modern diesel technologies;
they have not been matched to the inherited factors' implied fleet or
fuel-carbon conventions. Therefore
the contrast is a combined factor/fleet-configuration sensitivity, not an
isolated comparison of mathematical formulations or evidence of superior
accuracy. Parcel utilization for emissions remains the supplied constant 0.5, rather
than being recomputed from the decreasing parcel count. Consequently, this
experiment does not evaluate dynamic parcel payload emissions, and zero-
utilization activity is not a complete measure of empty parcel travel.
No non-exhaust factors are supplied. Nominal CO$_2$ mass is reported,
not CO$_2$ equivalents or a life-cycle footprint.

Missing inherited road-type factors otherwise yield zero modelled emissions.
The primary cross-configuration comparison therefore uses a fixed common
physical-road support, defined as the intersection of CO$_2$ road classes
covered for every assignment vehicle, excluding connectors and ferries.
Unrestricted totals and support coverage are retained as diagnostics.
CO and pollutant-specific quantities are not treated as interchangeable across
methods; in particular, the inherited PM boundary requires reconciliation
before comparison with COPERT exhaust PM. A flat-grade COPERT calculation on
the inherited, zero-penalty activity additionally separates direct factor
response to terrain from changes in selected paths.

\subsection{Experimental design and reproducibility}
Two schedulers, three path-cost coefficients and two emissions configurations
yield 12 configurations per upstream realization. Both emissions calculations
reuse exactly the same link activity, speed and payload representation.
Link-level vehicle-kilometres, rather than skim-based tour distances, define
network activity outcomes. Parcel and whole-network outcomes are retained
separately. Scheduling contrasts compare paired realizations at a fixed path
coefficient and emissions configuration; path-cost contrasts retain scheduler
and emissions configuration. Factor contrasts retain the same assignment.

These five realizations constitute an exploratory pilot. Means and observed
replicate ranges describe conditional simulation variability; they do not
establish stochastic equivalence to the single source realization or a
confirmatory precision target. Links and time intervals are not independent
replications. Confirmatory stochastic claims would require a prespecified precision
criterion and a justified replication plan.

To bound memory use, emissions are evaluated in row-local batches using the
frozen model functions, with the CLI's float32 output precision and float64
aggregation. Cached Dijkstra predecessors are keyed by the exact sparse cost
graph, ordering, numerical representation and SciPy version. Cached/uncached
assignment and batched/full-frame emissions are checked for numerical identity
on controlled fixtures. These execution adaptations do not change the model
formulation. Inputs, output hashes, commands, timings and environment versions
are retained with the experiment records.
